% Chapter fragment: input this file from the thesis; do not add a second document environment.
\section{Background and Related Work}
\label{sec:background}

New training data affect more than a continual learner's performance on its
current task. They can also change what it remembers from earlier tasks.
This makes the training data a possible route for an attack on knowledge
the model has already acquired. Detecting such an attack requires some
understanding of how the learner updates, what the attacker can change,
and which model responses are available before training begins.

\subsection{Continual Learning and Forgetting}
\label{subsec:bg-continual-learning}

Continual learning studies models that learn from a sequence of tasks or
data distributions. The difficulty is to learn new information while
retaining earlier knowledge. Training on a new task can change shared
parameters in ways that substantially reduce accuracy on previous tasks,
an effect known as catastrophic forgetting~\cite{parisi2019continual}.
Some protocols allow earlier examples to be stored and replayed; others
require learning to continue without them.

Let $D_t=\{(x_i,y_i)\}_{i=1}^{N_t}$ denote the training set of task $t$,
with tasks indexed from zero. A learner enters task $t$ with parameters
$\theta_{t-1}$ and produces $\theta_t$ after training on $D_t$. Minimising
the current empirical loss,
\begin{equation}
    \mathcal{L}_t(\theta)
    = \frac{1}{N_t}\sum_{(x,y)\in D_t}
      \ell\bigl(f_\theta(x,t),y\bigr),
    \label{eq:bg-current-loss}
\end{equation}
places no explicit constraint on performance on $D_0,\ldots,D_{t-1}$.
Here the argument $t$ allows a task-specific output head to be selected;
the class of $x$ remains to be predicted.

The balance between retaining earlier knowledge and adapting to new data
is usually described as the stability--plasticity trade-off. Replay methods
revisit stored or generated examples. Parameter-isolation methods allocate
model capacity to different tasks. Regularisation methods penalise changes
to parameters or predictions associated with earlier knowledge
\cite{parisi2019continual}. The experiments here focus on this last group
and examine its behaviour when the new training data are deliberately
modified.

\subsubsection{Task Identity and Class Labels}

Van de Ven and Tolias distinguish three evaluation scenarios according to the
information available at test time~\cite{vandeven2019scenarios}. In
task-incremental learning, task identity is provided and can select the
appropriate output head. In class-incremental learning, the model must
classify among the classes encountered so far without being told the task.
In domain-incremental learning, the input distribution changes while the
required output space remains shared, and task identity is not provided.
These scenarios impose different demands on the classifier and need to be
distinguished when comparing results.

The learner used here is task-incremental and has a separate head for each
task. The incoming task and its output space are known, but the detector
must score images without their true class labels. Its supervised training
uses a different label: whether an image is clean or artificially poisoned.
Those labels are available in the controlled source experiment and absent
when a new dataset is scored. Task identity, image class and clean/poison
status therefore have separate roles throughout the thesis.

\subsubsection{Measuring Retention}

Let $a_{t,j}$ be the test accuracy on task $j$ after learning task $t$.
The resulting accuracy matrix records both learning and retention. Following
the backward-transfer formulation of Lopez-Paz and Ranzato
\cite{lopezpaz2017gem}, the average change on earlier tasks after learning
task $t\geq 1$ is
\begin{equation}
    \operatorname{BWT}_t
    = \frac{1}{t}\sum_{j=0}^{t-1}(a_{t,j}-a_{j,j}).
    \label{eq:bg-bwt}
\end{equation}
A negative value means that earlier-task accuracies have declined, on
average, from their immediate post-training values. The mean historical
accuracy is
\begin{equation}
    A_t^{\mathrm{past}}=\frac{1}{t}\sum_{j=0}^{t-1}a_{t,j},
    \label{eq:bg-past-accuracy}
\end{equation}
which measures the performance that remains after the update. Both quantities
are useful: a small backward-transfer loss can hide poor historical accuracy
if the earlier tasks were never learned well.

Since clean sequential training can itself cause forgetting, poisoning needs
a paired control. Starting from the same pre-update checkpoint, the learner
is trained separately on clean and poisoned versions of the incoming task.
The difference
$A_{t,\mathrm{clean}}^{\mathrm{past}}-
 A_{t,\mathrm{poison}}^{\mathrm{past}}$ then describes the extra historical
accuracy loss attributable to the changed data in that experiment. Measuring
the benefit of a detector would require a further training run in which its
decisions are used to filter the data or modify learning.

\subsection{Regularisation-Based Continual Learning}
\label{subsec:bg-regularisation}

Several methods considered here use a parameter penalty of the general form
\begin{equation}
    \mathcal{L}_t^{\mathrm{CL}}(\theta)
    = \mathcal{L}_t(\theta)
      +\frac{\lambda}{2}\sum_r
       \Omega_{t-1,r}(\theta_r-\bar\theta_{t-1,r})^2.
    \label{eq:bg-regularisation}
\end{equation}
Here $r$ indexes scalar parameters, $\bar\theta_{t-1}$ is a retained reference
parameter vector, and $\Omega_{t-1,r}$ measures importance. Methods differ
in how these quantities are estimated and updated. A larger penalty restricts
movement in important coordinates, preserving earlier knowledge at the
possible cost of slower adaptation to the new task.

\subsubsection{Elastic Weight Consolidation}

Elastic Weight Consolidation (EWC) estimates parameter importance using a
diagonal approximation to the Fisher information~\cite{kirkpatrick2017ewc}.
For a previous task $A$, its two-task penalty is
\begin{equation}
    \mathcal{L}_{\mathrm{EWC}}(\theta)
    = \mathcal{L}_t(\theta)
      +\frac{\lambda}{2}\sum_r
       F_{A,r}(\theta_r-\theta_{A,r}^{*})^2.
    \label{eq:bg-ewc}
\end{equation}
Larger Fisher values give stronger constraints. The penalty follows from a
local approximation to the posterior after the earlier task. Using only the
diagonal makes it practical to compute, at the cost of omitting interactions
between parameters. Extending EWC to a task sequence also requires choosing
how to retain or combine importance estimates.

EWC's accumulated importance values control the training update. By
contrast, the detector's gradient norms describe the response of a frozen
model to one image. Although both involve gradients, they measure different
quantities.

\subsubsection{Memory Aware Synapses and Riemannian Walk}

Memory Aware Synapses (MAS) measures importance through the sensitivity of
the network's output function to parameter changes
\cite{aljundi2018mas}. Unlike a loss-based importance estimate, this
calculation can use unlabelled inputs. Parameters judged important are then
protected during subsequent learning. The unlabelled part of MAS is this
importance estimate; the task classifier is still trained with class labels.

Riemannian Walk (RWalk) combines Fisher-based information with a
training-path contribution to parameter importance
\cite{chaudhry2018rwalk}. Its formulation connects changes in model
distributions with the optimisation trajectory. Chaudhry et al. also
distinguish forgetting from intransigence, the difficulty of acquiring new
knowledge. A method that barely changes its parameters may retain earlier
predictions but learn the incoming task poorly, so retention and new-task
accuracy need to be examined together.

\subsubsection{Auxiliary-Network Approaches}

Active Forgetting with synaptic Expansion--Convergence (AFEC) addresses
negative transfer from earlier knowledge to a new task
\cite{wang2021afec}. It introduces additional parameters for learning the
new task and then selectively combines them with the existing network.
This allows earlier constraints to be relaxed where they obstruct learning
the new task.

Auxiliary Network Continual Learning (ANCL) also introduces an auxiliary
network, but uses it to provide a plasticity-oriented reference for the main
learner~\cite{kim2023ancl}. Its regulariser balances the previous model and
the auxiliary solution. The shared use of an auxiliary network thus serves
different formulations in AFEC and ANCL. Their EWC-based variants are
included in the experimental comparison.

\subsection{Data Poisoning and Forgetting-Inducing Attacks}
\label{subsec:bg-poisoning}

In a poisoning attack, the adversary changes data on which the victim will
train. In an evasion attack, the adversary instead changes an input to a
fixed, trained model. Poisoning therefore acts through the learner's update:
the modified data must cause damage after they have been used for training.
For example, Biggio et al. optimise an injected SVM training point to
increase a downstream validation loss~\cite{biggio2012poisoning}. An input's
difficulty before training is only part of this problem.

A poisoning objective can be written schematically as
\begin{align}
    \max_{\Delta\in\mathcal{C}}\quad&
      \mathcal{J}\bigl(\theta^{+}(\Delta)\bigr),\label{eq:bg-poison-outer}\\
    \text{where}\quad&
      \theta^{+}(\Delta)=
      \mathcal{A}(\theta^{-},D_t^{\Delta}).\label{eq:bg-poison-inner}
\end{align}
Here $\mathcal{A}$ is the victim's update procedure, $\theta^{-}$ is its
pre-update state, and $\mathcal{C}$ specifies the attacker's permitted
changes. The outer objective $\mathcal{J}$ measures the attacker's desired
damage. The update operator $\mathcal{A}$ can represent a finite training
schedule, as well as an approximate solution to a learning objective.

For additive image poisoning, the modified input can be represented as
$\widetilde{x}_i=\Pi_{[0,1]^d}(x_i+\delta_i)$, with
$\|\delta_i\|_\infty\leq\varepsilon$. The projection keeps pixel values in
range. The budget $\varepsilon$ bounds the change to each pixel. The
poisoning fraction specifies how many images are changed. These are separate
settings: small perturbations may be applied to every image, or larger
perturbations to only a few.

\subsubsection{BrainWash}

BrainWash targets the retention of earlier tasks rather than only the
accuracy of the incoming task~\cite{abbasi2024brainwash}. The attacker has
access to the learned model and incoming-task data, but not the original
historical training sets. It first synthesises historical inputs through
model inversion, then uses them to optimise perturbations that make a
subsequent update damage earlier-task performance.

In the reckless setting, the attacker seeks forgetting without requiring
good performance on the incoming task. In the cautious setting, the attack
also seeks to preserve incoming-task accuracy, so monitoring that accuracy
alone may miss the damage. Both use constrained input perturbations,
optimised through an approximation to a bilevel objective
\cite{abbasi2024brainwash}.

For detection, the relevant point is that damage to historical tasks may
remain hidden by acceptable performance on the new task. Historical retention
and incoming-task accuracy are therefore measured separately. The attack
budgets and optimisation settings used here are given in the experimental
setup.

\subsubsection{Poisoning Detection and Existing Defences}

Backdoor defences provide related but different detection settings. Spectral
Signatures identifies unusual directions in learned representations and uses
them to detect and remove poisoned examples~\cite{tran2018spectral}. This
work uses representations to reveal contamination that ordinary test
accuracy can overlook. Its setting is a classifier trained on backdoored
data, whereas the detector here examines a new task before learning it.
Whether comparable signals exist at this earlier stage must be tested.

A detection score also needs to be distinguished from a training defence.
The former identifies suspicious data; the latter acts on the learning
procedure, for example by filtering images or constraining updates. The
experiments build on the PROACT repository, which already contains
training-time defence components~\cite{proactsoftware}. The detector is a
separate addition. Assessing whether its alerts improve retention would
require using them during training, beyond measuring detection and attack
effectiveness separately.

\subsection{Model Inversion as Historical Reference Data}
\label{subsec:bg-inversion}

Model inversion synthesises inputs using information contained in a trained
model. DeepInversion optimises image inputs while keeping the network
fixed~\cite{yin2020deepinversion}. A classification objective encourages
the requested class, while image priors and BatchNorm feature-statistic
constraints guide synthesis. The stored BatchNorm statistics help constrain
the otherwise underdetermined search: many inputs can produce a desired
network output.

The resulting images are synthetic inputs, rather than recovered copies of
the original training set. Matching selected feature statistics still leaves
room for differences from real images. Smith et al. examine this difficulty
in data-free class-incremental learning~\cite{smith2021dreaming}: replaying
inverted images with standard distillation is insufficient on common
benchmarks, motivating changes to the learning procedure. Their problem is
different from detection, but it illustrates why the quality of a synthetic
reference needs to be examined rather than assumed.

Inversion serves two purposes in these experiments. It supplies the attacker
with a surrogate for inaccessible historical data, and supplies the
label-free detector with historical reference images. These uses place
different demands on the samples. Images that help construct a damaging
update need not be representative enough to calibrate a clean-data test.
The reference-based detector consequently relies both on the historical
checkpoints being clean and on the adequacy of the synthetic references.

\subsection{Model-Based Descriptors and Detection Scores}
\label{subsec:bg-descriptors}

The frozen learner provides several ways to describe an incoming image.
Its output probabilities give the head's response, activations give its
internal representation, and loss gradients measure local parameter
sensitivity. These descriptors can be supplied to a classifier or compared
across datasets without training the learner on the incoming images.

\subsubsection{Gradients and Their Limits}

For a differentiable loss, write
$g(x,y)=\nabla_\theta\ell(f_\theta(x),y)$. A small gradient step on $(x,y)$
changes parameters approximately by $-\eta g(x,y)$. If
$\mathcal{L}_{\mathrm{past}}$ is a differentiable historical loss, a
first-order expansion gives
\begin{equation}
 \mathcal{L}_{\mathrm{past}}(\theta-\eta g)
 -\mathcal{L}_{\mathrm{past}}(\theta)
 \approx -\eta\,\nabla_\theta\mathcal{L}_{\mathrm{past}}(\theta)^\top g.
 \label{eq:bg-gradient-interference}
\end{equation}
The inner product determines whether this small step increases or decreases
historical loss. Gradient magnitude alone omits that direction information.
This is a local explanation of interference; a full training update also
involves successive steps, regularisation and changes in the representation.

Influence functions provide a related approach to studying how training
examples affect model predictions~\cite{koh2017influence}. They account
for local curvature through the training objective, rather than assigning
influence from a raw gradient norm alone. The descriptors used here retain
less information: they summarise gradients without estimating the curvature
or the eventual effect of an image on training.

Gradient magnitude can be summarised at different scales. If the backbone parameters are
partitioned into tensors $\theta^{(1)},\ldots,\theta^{(q)}$, one can record
$\|\nabla_{\theta^{(j)}}\ell\|_2$ for every tensor. This is one feature for
an entire weight or bias tensor, not one feature per scalar parameter. A
stage norm instead combines several tensors. For a disjoint partition,
\begin{equation}
 \|\nabla_\theta\ell\|_2^2
 =\sum_{j=1}^{q}\|\nabla_{\theta^{(j)}}\ell\|_2^2.
 \label{eq:bg-gradient-partition}
\end{equation}
Tensor norms retain the distribution of gradient magnitude within a stage,
which the stage norm combines into a single value. They also increase
dimensionality and depend on tensor size and parameterisation. Comparing
the two representations tests whether that additional detail helps detection.

Computing a loss gradient normally requires a class target. Since the
detector lacks the true incoming class, it uses the model's predicted class
as a fixed pseudo-target for differentiation. The gradient then measures
sensitivity relative to that prediction, including when it is incorrect.
The method chapter specifies this construction and the fixed incoming head
used to obtain it.

\begin{samepage}
\subsubsection{Predictions and Activations}

For predicted probabilities $p$, entropy is
\mbox{$H(p)=-\sum_c p_c\log p_c$}. Other summaries are the maximum probability
and the gap between the two largest values. Hendrycks and Gimpel use maximum
softmax probability as a baseline for identifying misclassified and
out-of-distribution inputs~\cite{hendrycks2017baseline}. Low confidence
can thus be useful for identifying unfamiliar inputs, although unfamiliarity
can arise from clean task changes as well as poisoning.
\par
\end{samepage}

Guo et al.~\cite{guo2017calibration} show that modern trained neural networks
can produce probabilities that are poorly calibrated to classification
correctness. Here the incoming head has not yet been trained on its task,
so its probability vector is used as a model response rather than a
calibrated estimate of correctness. The activation norm adds a measurement
of representation magnitude. Any association between these responses and
poisoning is learned or tested in the detection experiment.

\subsubsection{Ranking, Calibration and Dataset Decisions}

A supervised detector maps descriptors $z(x)$ to a scalar score, for
example through logistic regression,
\begin{equation}
 s(x)=\sigma\bigl(w^\top\widetilde z(x)+b\bigr),\qquad
 \sigma(u)=\frac{1}{1+e^{-u}},
 \label{eq:bg-logistic-score}
\end{equation}
where $\widetilde z$ denotes descriptors transformed using source-fitted
scaling. The binary target distinguishes clean and poisoned images. The
score lies between zero and one, but its probabilistic interpretation depends
on the training population and class proportions. A balanced source
experiment alone cannot calibrate it for deployment on a different task.

ROC analysis evaluates how examples are ordered as the decision threshold
varies~\cite{fawcett2006roc}. A high ROC-AUC can therefore coexist with
poor decisions at one frozen cutoff. After task transfer, for example, clean
and poisoned scores may both shift while their relative ordering remains
largely unchanged.

Probability calibration concerns agreement between predicted probabilities
and observed frequencies~\cite{guo2017calibration}. Threshold calibration
has a narrower purpose: choosing a cutoff for a decision rule. Here that
cutoff is estimated from historical data and frozen before the target task.
Its error rate on the target remains something to measure.

A sample score can be turned into a dataset alert by counting the images
above a cutoff $u$:
\begin{equation}
 K(D)=\sum_{i=1}^{n}\mathbf{1}\{s(x_i)\geq u\}.
 \label{eq:bg-dataset-count}
\end{equation}
A count rule alerts when $K(D)\geq k$. Under independent clean-image
alarms with a common probability $p$, the count has a
$\operatorname{Binomial}(n,p)$ distribution. A one-sided Clopper--Pearson
bound can account for uncertainty in $p$ estimated from a separate clean
sample~\cite{clopper1934binomial}. This calculation relies on the clean-image
alarm probability remaining stable. Task shift can violate that assumption.
The number of distinct calibration images also matters: repeatedly drawing
datasets from the same finite pool does not supply new independent images.

\subsection{Historical-Reference Distribution Tests}
\label{subsec:bg-reference-tests}

Without labelled clean/poison examples, historical data offer a possible
reference for incoming descriptors. Such a comparison tests a property of
the data distributions, rather than the cause of a difference. New classes
and imperfect inversion samples can produce clean distribution shifts;
conversely, poisoning may leave the property being tested almost unchanged.

\subsubsection{Maximum Mean Discrepancy}

Maximum mean discrepancy (MMD) compares distributions through their
embeddings in a reproducing kernel Hilbert space
\cite{gretton2012mmd}. For independent $Z,Z'\sim P$ and $W,W'\sim Q$,
its squared population value is
\begin{equation}
 \operatorname{MMD}^2(P,Q)
 =\mathbb{E}k(Z,Z')+\mathbb{E}k(W,W')-2\mathbb{E}k(Z,W).
 \label{eq:bg-mmd}
\end{equation}
A characteristic kernel distinguishes unequal distributions at the
population level. A finite-sample test still depends on sample size,
kernel choice and the reference data. Random Fourier features approximate
shift-invariant kernels by finite-dimensional inner products, reducing
the cost of repeated comparisons~\cite{rahimi2007random}.

This sensitivity also presents a difficulty for historical references.
Synthetic historical images and real incoming images may differ even when
both are clean. In that case, rejecting equality is consistent with the
statistical question but unhelpful for identifying poisoning. A permutation
test requires exchangeability under its null, which the clean status of
the two groups alone cannot establish. Clean-alert and attack-detection
rates are therefore both needed to assess this baseline.

\subsubsection{Kendall Correlations and Rank Dependence}

Kendall's rank correlation compares the ordering of paired observations
\cite{kendall1938rank}. For descriptor coordinates $a$ and $b$, the
tau-a statistic is
\begin{equation}
 \widehat\tau_{ab}
 =\frac{2}{n(n-1)}\sum_{i<j}
   \operatorname{sign}(z_{ia}-z_{ja})
   \operatorname{sign}(z_{ib}-z_{jb}).
 \label{eq:bg-kendall}
\end{equation}
Concordant orderings contribute $+1$, discordant orderings $-1$, and ties
zero. Tau-a is not the tie-normalised tau-b statistic. A vector of these
coefficients describes pairwise rank dependence across a dataset.

Strictly increasing transformations of individual coordinates leave these
orderings unchanged. A rank comparison ignores these marginal changes,
allowing some scale differences between tasks to be removed. It also misses
an attack expressed entirely through such transformations. More generally,
pairwise Kendall coefficients describe only part of the joint distribution;
higher-order dependence may change while these coefficients stay fixed.

I use the name \emph{Rank} for the historical-reference test built from
these coefficients. The descriptors are measured per image, but the
correlations summarise a whole dataset. This gives a dataset-level
comparison without first classifying individual images as clean or
poisoned. The test statistic and the rule for combining historical
comparisons are described in the method chapter.

\subsubsection{Jackknife Variation and Bootstrap Approximation}

Correlation estimates vary with the sampled images. The jackknife measures
this variation by leaving out observations in turn
\cite{efron1981jackknife}. For a statistic $T_n$
and its leave-one-out value $T_{n,-i}$, the pseudovalue is
\begin{equation}
 P_i=nT_n-(n-1)T_{n,-i}.
 \label{eq:bg-jackknife}
\end{equation}
Variation across the pseudovalues estimates first-order variability in the
statistic. When that variation is effectively zero, a first-order
approximation becomes degenerate and needs separate treatment.

Gaussian multiplier methods approximate sampling errors using random
weights applied to estimated observation contributions. Chen develops
Gaussian and bootstrap approximations for nondegenerate, high-dimensional
U-statistics, including jackknife-based multiplier procedures
\cite{chen2018bootstrap}. Kendall coefficients are order-two U-statistics,
which makes this framework relevant to comparing correlation profiles.
Applying it here still requires assumptions about independence and
nondegeneracy. These are particularly important for synthetic inversion
samples. The implemented two-sample procedure is an approximation, rather
than a finite-sample guarantee supplied by Chen's result.

The bootstrap draws simulate fluctuations of the statistic conditional on
the observed images; the image sample itself stays fixed. The resulting
p-value concerns a null hypothesis about rank dependence, not the
probability of poisoning. Rejection and non-rejection must be interpreted
within that scope, with non-rejection leaving open changes that the test
has little power to detect.

\subsection{Scope of This Thesis}
\label{subsec:bg-scope}

The remaining question is how useful these measurements are before a new
task is learned. The supervised approach asks whether source-trained scores
and historical thresholds transfer to a later task. The label-free approach
asks whether inversion references support useful dataset comparisons without
examples labelled as poisoned. Both depend on how model responses change
across tasks, but they use different information and make different assumptions.

The CL algorithms, BrainWash attack and inversion procedure are existing
components. This work develops the detector around them, compares gradient
granularities and feature groups, and examines where the resulting decisions
work or fail. The next chapter defines the information available to the
detector and attacker before describing these methods in detail.
